% TOP_PROBLEM: {"id": 1496, "title": "Odd Perfect Numbers", "category_id": 1, "category": {"id": 1, "name": "number_theory", "display_name": "Number Theory", "description": "Properties of integers, prime numbers, Diophantine equations.", "slug": "number-theory", "order_index": 1, "created_at": "2026-07-31T15:26:25.670Z"}, "status": "open"}
% FIRST_POSTED: null
% ENTRY_KIND: statement_only
% Imported from: batch08.tex; UnsolvedMath ID 1496
\documentclass[11pt]{article}
\usepackage[margin=25mm]{geometry}
\usepackage{fontspec}
\setmainfont{Latin Modern Roman}
\usepackage{amsmath,amssymb,mathtools}
\usepackage{unicode-math}
\setmathfont{Latin Modern Math}
\usepackage{xurl,hyperref}
\hypersetup{colorlinks=true,urlcolor=blue}
\setcounter{secnumdepth}{0}
\setlength{\parindent}{0pt}
\setlength{\parskip}{5pt}
\setlength{\emergencystretch}{3em}
\newcommand{\sref}[2]{\hyperlink{source:#1:#2}{[S#2]}}
\newcommand{\cataloguescope}[1]{\paragraph{Recorded target.}#1}
\begin{document}
% UnsolvedMath ID 1496; source catalogue code NUM-002
\setcounter{section}{1495}
\section{Odd Perfect Numbers}\label{1496}
{\small\sffamily UnsolvedMath ID 1496 · Number Theory}
\subsection{Definitions and mathematical statement}

\paragraph{Shared notation.}$\mathbb N=\{1,2,\ldots\}$, and $\mathbb Z,\mathbb Q,\mathbb R,\mathbb C$ denote the integers, rationals, reals and complex numbers. Any other conventions are specified locally.


For \(n\in\mathbb N\), let
\[
 D(n)=\{d\in\mathbb N:d\mid n\},\qquad
 \sigma(n)=\sum_{d\in D(n)}d.
\]
The proper positive divisors are \(D(n)\setminus\{n\}\).
A perfect number satisfies \(\sum_{d\mid n,\,d<n}d=n\), equivalently
\(\sigma(n)=2n\).
\paragraph{Statement recorded in the supplied source.}
Does there exist an odd integer \(n\ge1\) such that
\[
                            \sigma(n)=2n?
\] \sref{1496}{2} \sref{1496}{3}
\subsection{Short English statement}
Is any odd positive integer equal to the sum of its proper positive divisors?
\subsection{Sources}
\begin{description}
\item[\hypertarget{source:1496:1}{S1}] ULAM.AI and the UnsolvedMath Contributors, \emph{UnsolvedMath: A Curated Collection of Open Mathematics Problems}, record 1496 (NUM-002). CC BY 4.0. \url{https://huggingface.co/datasets/ulamai/UnsolvedMath}.
\item[\hypertarget{source:1496:2}{S2}] Pascal Ochem and Joshua Zelinsky, \emph{On odd perfect numbers with exactly one even exponent greater than 2}, arXiv:2607.19746 (2026), Section 1, definition of an odd perfect number. The cited definition concerns all odd perfect numbers, independently of the paper's restricted factorization theorem. \url{https://arxiv.org/abs/2607.19746}.
\item[\hypertarget{source:1496:3}{S3}] John Cowles and Ruben Gamboa, \emph{Perfect Numbers in ACL2}, EPTCS 192 (2015), 53--59, DOI 10.4204/EPTCS.192.5; arXiv:1509.06081v1, Sections 1--2, the unrestricted infinitude question and $\sigma(n)=2n$ definition. \url{https://arxiv.org/abs/1509.06081v1}.
\end{description}
\label{end:1496}
\end{document}
